\section{Introduction}
A rare favorable report has no operational meaning until the experiment that obtains it has been paid for. A predictive state has no finite-memory meaning until its update, legal commands, stopping decision and task readout have been implemented. These two observations lead to the same question: which geometry of an actually prepared causal experiment controls its attainable prediction risk under simultaneous acquisition and memory constraints?

There is a further obstruction when the hidden dynamics are unknown. A lower bound on conditional audit variance can fail after a highly informative history, even though a single program remains uniformly minimax. Imposing a positive variance floor on every history can therefore confuse a sufficient pointwise calibration mechanism with the resource law itself. This paper separates the two. The quantities to balance are the report measure and the \emph{unnormalized} squared-loss measure. Balancing the first without the second does not survive report-dependent selection.

The central result, Theorem~\ref{thm:main}, starts with positive prepared kernels, a declared scored future test, paid calls and a physical stopping interface. It derives a density-limited acquired measure from a fixed admissible reference model. It then constructs one common rejection program from a task-balancing identity certified before choosing the memory budget. In the exact case the checkpoint risk is $\Qm_m(N/\kappa)$ and the autonomous risk is $\Qm_{M-B}(N/\kappa)$, where $B$ is a proved number of mandatory live states. The reference model is a genuine member of the uncertainty class, not an auxiliary distribution selected after seeing the code. Conversely, the implemented upper does not know that model. Its full hidden state may remain nonstationary and nonmixing.

The theorem also permits a finite family of report and readout transformations. Their exported terminal dictionary is counted: $m$ virtual values may require $Dm$ physical values. Under acquired covering and small-ball hypotheses this yields matching resource orders; in invariant examples the alphabet expansion can be avoided. The predictive quotient is constructed only after specifying executable tests and measurable realization conditions. A task statistic is not silently promoted to the full predictive quotient.

The principal additions are the common raw task-measure criterion, its exact stopped minimax completion and its verification without a positivity floor. Proposition~\ref{prop:moment} gives a necessary and sufficient linear criterion for a specified finite family of balancing commands; it is not a characterization of all possible causal observers. Proposition~\ref{prop:instrument} proves the integrated continuity needed at singular reports. Theorem~\ref{thm:classical} and Theorem~\ref{thm:quantum} verify the central theorem by genuinely different raw instruments. Theorem~\ref{thm:gaussian} strengthens the retained Gaussian application \cite{GTF52} to a zero-prior-floor uncertainty class and a ratio-one risk asymptotic: one terminal value cannot exploit waiting in the worst case, whereas two values achieve the asymptotically optimal waiting law, uniformly over arbitrary additional memory.

The underlying tools have clear antecedents. Randomization over a finite symmetry group is standard invariant decision theory \cite{Luschgy}; Pauli randomization is classical quantum information \cite{AMTW}. Neither is claimed as a new operation. Optimal rejection rules go back to Chow \cite{Chow}; thresholding conditional variance in square-loss regression with a fixed rejection rate is established in \cite{DHZ}. Our Gaussian calculation does not claim either principle. Its additional assertions concern predictable paid waiting, a common rule over history-dependent hidden dynamics, the exact one-value obstruction and an asymptotic sharpness statement uniform over all larger finite terminal alphabets. Occupation-measure Markovization and absorbing control are developed in \cite{PZExtreme,PZContinuity,DPAbsorb,DPChatter}. The present proof does not identify all hidden histories with a known finite driver or infer compactness of a general occupation set. It uses predictable counts under one witness and history-uniform task-measure identities for one common upper program. General finite-memory POMDP reachability remains a different and computationally difficult problem \cite{CCGK,ACSS}. The marked reliability--resolution calculation and physical-port accounting were developed in the retained causal-acquisition and predictable-geometry manuscripts \cite{GTF52,GTF54}; their proofs and original scope are retained and the needed arguments are repeated here.

The limits are also mathematical. Existence of a balancing command and of a suitable reference member are substantive assumptions. No universal learning rate for unknown kernels, no unrestricted equivalence of interventions on different targets, and no efficient synthesis theorem is asserted. The exact profile is first a theorem for Borel finite-state transducers. Proposition~\ref{prop:finite} supplies finite-description implementations under explicit effective report and probability certificates, with their precision and workspace charged separately. This distinction is necessary even for ordinary continuous-observation experiments.
